%%
%% This is file `sample-sigconf.tex',
%% generated with the docstrip utility.
%%
%% The original source files were:
%%
%% samples.dtx  (with options: `all,proceedings,bibtex,sigconf')
%% 
%% IMPORTANT NOTICE:
%% 
%% For the copyright see the source file.
%% 
%% Any modified versions of this file must be renamed
%% with new filenames distinct from sample-sigconf.tex.
%% 
%% For distribution of the original source see the terms
%% for copying and modification in the file samples.dtx.
%% 
%% This generated file may be distributed as long as the
%% original source files, as listed above, are part of the
%% same distribution. (The sources need not necessarily be
%% in the same archive or directory.)
%%
%%
%% Commands for TeXCount
%TC:macro \cite [option:text,text]
%TC:macro \citep [option:text,text]
%TC:macro \citet [option:text,text]
%TC:envir table 0 1
%TC:envir table* 0 1
%TC:envir tabular [ignore] word
%TC:envir displaymath 0 word
%TC:envir math 0 word
%TC:envir comment 0 0
%%
%% The first command in your LaTeX source must be the \documentclass
%% command.
%%
%% For submission and review of your manuscript please change the
%% command to \documentclass[manuscript, screen, review]{acmart}.
%%
%% When submitting camera ready or to TAPS, please change the command
%% to \documentclass[sigconf]{acmart} or whichever template is required
%% for your publication.
%%
%%
\documentclass[sigplan,anonymous,review,nonacm]{acmart}
%%
%% \BibTeX command to typeset BibTeX logo in the docs
\AtBeginDocument{%
  \providecommand\BibTeX{{%
    Bib\TeX}}}

%% Rights management information.  This information is sent to you
%% when you complete the rights form.  These commands have SAMPLE
%% values in them; it is your responsibility as an author to replace
%% the commands and values with those provided to you when you
%% complete the rights form.
\setcopyright{acmlicensed}
\copyrightyear{2018}
\acmYear{2018}
\acmDOI{XXXXXXX.XXXXXXX}
%% These commands are for a PROCEEDINGS abstract or paper.
\acmConference[Conference acronym 'XX]{Make sure to enter the correct
  conference title from your rights confirmation email}{June 03--05,
  2018}{Woodstock, NY}
%%
%%  Uncomment \acmBooktitle if the title of the proceedings is different
%%  from ``Proceedings of ...''!
%%
%%\acmBooktitle{Woodstock '18: ACM Symposium on Neural Gaze Detection,
%%  June 03--05, 2018, Woodstock, NY}
\acmISBN{978-1-4503-XXXX-X/2018/06}


%%
%% Submission ID.
%% Use this when submitting an article to a sponsored event. You'll
%% receive a unique submission ID from the organizers
%% of the event, and this ID should be used as the parameter to this command.
%%\acmSubmissionID{123-A56-BU3}

%%
%% For managing citations, it is recommended to use bibliography
%% files in BibTeX format.
%%
%% You can then either use BibTeX with the ACM-Reference-Format style,
%% or BibLaTeX with the acmnumeric or acmauthoryear sytles, that include
%% support for advanced citation of software artefact from the
%% biblatex-software package, also separately available on CTAN.
%%
%% Look at the sample-*-biblatex.tex files for templates showcasing
%% the biblatex styles.
%%

%%
%% The majority of ACM publications use numbered citations and
%% references.  The command \citestyle{authoryear} switches to the
%% "author year" style.
%%
%% If you are preparing content for an event
%% sponsored by ACM SIGGRAPH, you must use the "author year" style of
%% citations and references.
%% Uncommenting
%% the next command will enable that style.
%%\citestyle{acmauthoryear}


%%
%% end of the preamble, start of the body of the document source.
\begin{document}

%%
%% The "title" command has an optional parameter,
%% allowing the author to define a "short title" to be used in page headers.
\title{COLoRA}

%%
%% The "author" command and its associated commands are used to define
%% the authors and their affiliations.
%% Of note is the shared affiliation of the first two authors, and the
%% "authornote" and "authornotemark" commands
%% used to denote shared contribution to the research.
% \author{Ben Trovato}
% \authornote{Both authors contributed equally to this research.}
% \email{trovato@corporation.com}
% \orcid{1234-5678-9012}
% \author{G.K.M. Tobin}
% \authornotemark[1]
% \email{webmaster@marysville-ohio.com}
% \affiliation{%
%   \institution{Institute for Clarity in Documentation}
%   \city{Dublin}
%   \state{Ohio}
%   \country{USA}
% }
\author{Anonymous Submission}

%%
%% By default, the full list of authors will be used in the page
%% headers. Often, this list is too long, and will overlap
%% other information printed in the page headers. This command allows
%% the author to define a more concise list
%% of authors' names for this purpose.
% \renewcommand{\shortauthors}{Trovato et al.}

%%
%% The abstract is a short summary of the work to be presented in the
%% article.
\begin{abstract}
  TODO
\end{abstract}

%%
%% The code below is generated by the tool at http://dl.acm.org/ccs.cfm.
%% Please copy and paste the code instead of the example below.
%%
\begin{CCSXML}
<ccs2012>
   <concept>
       <concept_id>10010520.10010521.10010542.10010546</concept_id>
       <concept_desc>Computer systems organization~Heterogeneous (hybrid) systems</concept_desc>
       <concept_significance>500</concept_significance>
       </concept>
   <concept>
       <concept_id>10011007.10010940.10010941.10010949.10010950</concept_id>
       <concept_desc>Software and its engineering~Memory management</concept_desc>
       <concept_significance>500</concept_significance>
       </concept>
   <concept>
       <concept_id>10011007.10010940.10010941.10010949.10010957.10010688</concept_id>
       <concept_desc>Software and its engineering~Scheduling</concept_desc>
       <concept_significance>500</concept_significance>
       </concept>
 </ccs2012>
\end{CCSXML}

\ccsdesc[500]{Computer systems organization~Heterogeneous (hybrid) systems}
\ccsdesc[500]{Software and its engineering~Memory management}
\ccsdesc[500]{Software and its engineering~Scheduling}

%%
%% Keywords. The author(s) should pick words that accurately describe
%% the work being presented. Separate the keywords with commas.
\keywords{TODO}
%% A "teaser" image appears between the author and affiliation
%% information and the body of the document, and typically spans the
%% page.
% \begin{teaserfigure}
%   \includegraphics[width=\textwidth]{sampleteaser}
%   \caption{Seattle Mariners at Spring Training, 2010.}
%   \Description{Enjoying the baseball game from the third-base
%   seats. Ichiro Suzuki preparing to bat.}
%   \label{fig:teaser}
% \end{teaserfigure}

\received{20 February 2007}
\received[revised]{12 March 2009}
\received[accepted]{5 June 2009}

%%
%% This command processes the author and affiliation and title
%% information and builds the first part of the formatted document.
\maketitle

\section{Introduction}
\label{sec:introduction}
Large-scale LLM serving increasingly combines two trends that are independently effective but jointly difficult to support: sparse Mixture-of-Experts (MoE) base models~\cite{todo:moe-llm} and LoRA-based adaptation~\cite{todo:lora-original,todo:lora-serving}. In production-style prefill--decode (PD) deployments~\cite{todo:pd-serving}, this combination is especially painful on the decode worker. Decode is latency-sensitive, runs with much smaller effective batches than prefill, and repeats the same routing and parameter-management decisions once per generated token. As a result, control-path inefficiencies that are tolerable in prefill can directly inflate token time and, more importantly, P99 tail latency in decode.

Existing serving systems do not fully capture this regime because they are built around the wrong managed object. In pure MoE serving, the natural unit is the expert. In pure multi-LoRA serving, the natural unit is the adapter. Existing multi-LoRA systems such as S-LoRA and VaLoRA~\cite{todo:slora,todo:valora} are therefore largely organized around adapter-centric management. In MoE plus multi-LoRA decode serving, however, the request only needs the LoRA residual associated with the specific experts chosen by the router. The effective unit of reuse is therefore a joint expert--adapter object rather than an expert-only or adapter-only object. This joint key expansion enlarges the working set, fragments reuse, and makes GPU residency much harder to maintain under a fixed memory budget.

Our trace-driven case study shows that this shift is not a small perturbation. When the access stream is modeled with joint keys, the top 1000 objects capture only 40.7\% of accesses under an independent adapter mapping and 32.6\% under a correlated mapping, versus 89.1\% for expert-only keys. The same trend appears in temporal locality: the finite P99 reuse distance rises from 1,393 events in the expert-only stream to 2,041 and 4,521 events in the two joint-key settings, while cold first touches increase from 1.0 to 9.7 and 13.6 per 10k events. These numbers show that multi-LoRA does not merely add more parameters. It pushes the decode worker into a more fragmented locality regime.

This fragmentation quickly becomes a systems problem under realistic cache budgets. Under a shared 2048-object LRU budget, cache replay yields a miss rate of 0.00159 for expert-only objects, but 0.01095 and 0.01476 for the two joint-key settings, which is 6.9$\times$ and 9.3$\times$ worse. A controlled synthetic sweep further shows that the tail penalty turns on early: at the 8192-object slice, the joint P99 latency-proxy penalty rises from 0.0\,ms at one LoRA to 411.7\,ms at four LoRAs, then stays on a broad plateau around 429.4\,ms for larger adapter counts. At the request level, the tail is dominated by cold-object touches rather than mild average degradation: at the 65536-object cold-miss floor, an average request touches 83.6 and 116.8 cold joint objects, while a P95+ tail request touches 912.8 and 1154.0. Together, these results suggest that misses are not accidental outliers in this regime. They are structural events concentrated on unlucky tail requests.

Once misses become structural, the usual load-then-run assumption becomes questionable. GPU-only systems that treat host memory primarily as spill space or a prefetch source~\cite{todo:gpu-offload} implicitly assume that a miss should be resolved by promoting weights back to GPU before execution. That assumption is often reasonable for reusable hot objects. It is far less attractive for long-tail decode misses, where the request batch is small, the object is cold, the reuse horizon is uncertain, and the fixed cost of host--device transfer is exposed directly on the critical path. The bottleneck therefore shifts from miss avoidance alone to miss handling itself.

This paper proposes COLoRA, a decode-worker runtime that treats a miss as a runtime execution choice instead of a mandatory weight-migration event. COLoRA maintains a bounded GPU hot cache for reusable joint objects and a CPU-resident full replica for all joint objects. On a hit, the request follows the fast GPU path. On a miss, the request does not block waiting for promotion by default. Instead, COLoRA executes the LoRA residual on the CPU while the base expert path remains on the GPU, and it promotes the missed object asynchronously only when its recent accesses indicate that promotion is likely to pay off. The result is a dual-path execution model that aims to absorb cold misses with a more stable fallback path while preserving GPU efficiency for the hot set.

The paper makes three contributions:
\begin{itemize}
\item We identify the joint-object fragmentation regime in MoE multi-LoRA decode serving and show, through a trace-driven case study, that the main failure mode is tail amplification from structural misses rather than a uniform slowdown.
\item We present COLoRA, a runtime built around asymmetric joint-object residency, CPU-first non-blocking miss handling, and selective asynchronous promotion.
\item We outline an evaluation methodology centered on tail latency, miss-handling breakdown, and overlap efficiency so that the paper can argue not only that COLoRA helps, but also why it helps and where it stops helping.
\end{itemize}


\section{Background and Motivation}
\label{sec:background}

\subsection{Serving Setting and System Boundary}
\label{subsec:setting}

This paper focuses on a single decode worker in a PD-separated serving stack. The upstream scheduler has already assigned requests to this worker, and the worker is responsible for executing decode tokens for an MoE base model while honoring per-request LoRA adapters. The design problem is therefore local to one worker: how should it coordinate GPU memory, CPU memory, GPU compute, and CPU compute so that decode remains stable under a fragmented access stream?

This boundary matters because it keeps the paper honest. COLoRA does not solve cluster-wide routing, admission control, or global KV placement. It also does not attempt to replace GPU execution for the dense backbone. The base MoE computation still runs on the GPU. The paper instead studies the part of the system that becomes unstable once MoE routing and multi-LoRA adaptation interact inside a decode worker: the placement and execution of the expert-specific LoRA residuals.

\subsection{Why the Managed Object Becomes a Joint Expert--Adapter Object}
\label{subsec:joint-object}

The first background point is conceptual. In a dense model with a single adapter, the adapter is a reasonable management unit because most tokens see the same adapted weights. In an MoE model, however, each token only activates a small routed subset of experts at each layer. Once multiple adapters are active at the same time, the worker does not need an entire adapter or an entire expert uniformly. It needs the LoRA residual associated with the specific expert selected for the specific request. That is the logical joint expert--adapter object that COLoRA is designed around.

This distinction is easy to miss but crucial to the paper story. Adapter-only management is too coarse because it keeps many cold expert slices resident together with a small hot subset. Expert-only management is also inaccurate because it ignores the adapter dimension and therefore overestimates reuse. The logical conclusion is that the system should reason about joint objects. In the implementation, this logical unit is materialized as projection-specific cache entries keyed by $(\mathit{projection}, \mathit{adapter}, \mathit{layer}, \mathit{expert})$, which lets COLoRA preserve the paper-level joint-object abstraction while matching how gate, up, and down LoRA weights are actually stored and executed.

\subsection{Trace-Driven Case Study: Joint Objects Fragment Decode Locality}
\label{subsec:case-study}

The case study exists to establish the failure mode before introducing the mechanism. We replay one aligned request stream under three object definitions: \texttt{B0} uses expert-only objects $(\mathit{layer}, \mathit{expert})$, \texttt{B1} uses joint objects $(\mathit{layer}, \mathit{expert}, \mathit{adapter})$ under an independent adapter mapping, and \texttt{B2} uses the same joint key under a correlated mapping. This setup lets the paper isolate what changes when adapters enter the access key while keeping the router trace fixed.

The first finding is that joint keying flattens popularity much earlier than an expert-only view would suggest. In \texttt{B0}, the hot head of the object distribution remains dominant, with 89.1\% of accesses absorbed by the top 1000 objects. In \texttt{B1} and \texttt{B2}, that coverage drops to 40.7\% and 32.6\%. The message is not simply that there are more objects. The stronger point is that the hot set itself becomes less concentrated, which immediately reduces the effectiveness of any bounded GPU cache.

The second finding is that temporal locality deteriorates in the tail even when average request structure does not look radically different. Under joint keys, finite reuse becomes much farther away and cold first touches become much more common. The finite P99 reuse distance rises from 1,393 events in \texttt{B0} to 2,041 in \texttt{B1} and 4,521 in \texttt{B2}, while cold first touches grow from 1.0 to 9.7 and 13.6 per 10k events. This result matters because decode latency is tail-sensitive: a few long reuse distances can be more damaging than a modest shift in the median.

The third finding is that locality loss translates directly into miss amplification and tail concentration under a shared cache budget. At 2048 objects, cache replay yields miss rates of 0.00159 for \texttt{B0}, 0.01095 for \texttt{B1}, and 0.01476 for \texttt{B2}. In the synthetic adapter-count sweep, the P99 latency-proxy penalty appears quickly once the workload leaves the single-LoRA regime and then stays high instead of decaying away with more cardinality. At the request level, the tail requests touch an order of magnitude more cold joint objects than the average request. These observations motivate the central design question of the paper: if a decode worker cannot reliably avoid misses in the joint-object regime, how should it handle those misses without making the tail worse?

\subsection{Takeaway for System Design}
\label{subsec:takeaway}

The background section leads to one concrete takeaway. In MoE multi-LoRA decode serving, the decode worker should not be optimized under the assumption that sufficiently aggressive caching will eliminate most important misses. The more realistic target is to keep the reusable hot set on GPU while making cold misses survivable. That target naturally suggests a heterogeneous design in which miss handling becomes an execution-path problem, not just a memory-placement problem.



\section{Challenges}
\label{sec:challenges}

The above observations define a new systems problem, but turning that problem statement into a useful runtime still requires solving three technical challenges.

\subsection{Challenge 1: Fine-Grained Residency Without Unstable Churn}
\label{subsec:challenge-residency}

The first challenge is to manage residency at the right granularity without making the cache thrash itself to death. Joint-object management is necessary for precision, but finer granularity also means more distinct keys, more misses, and more opportunities to waste bandwidth on objects that never become hot again. A practical design therefore needs to preserve joint-object precision while filtering out low-value promotions and suppressing repeated requeueing of the same cold object. Otherwise, the control path degenerates into churn rather than useful adaptation.

\subsection{Challenge 2: Making CPU Fallback Viable at Decode Granularity}
\label{subsec:challenge-cpu}

The second challenge is that ``just run the miss on CPU'' is not automatically a good idea. Decode misses are small, irregular, and latency-sensitive. Their batch sizes are often too small to amortize framework overhead, tensor reshaping, or generic BLAS startup costs. If the fallback path is not explicitly optimized for this regime, CPU execution simply replaces one form of tail latency with another. The system therefore needs a CPU miss path whose fixed cost is low enough to be competitive for cold, low-reuse objects.

\subsection{Challenge 3: Overlapping Two Execution Paths Without Creating a New Tail}
\label{subsec:challenge-overlap}

The third challenge is synchronization. A dual-path runtime introduces new interactions between GPU execution, CPU execution, and host--device transfers. If the GPU has to stall for the CPU at every miss, the design collapses back into a serialized pipeline. If the CPU queue becomes overloaded, fallback itself becomes the new tail bottleneck. The runtime must therefore overlap GPU hits with CPU misses when dependencies allow it, while exposing enough queueing and overlap state to keep the tail understandable and controllable.

\section{Design Overview}
\label{sec:design-overview}

\subsection{Core Idea}
\label{subsec:core-idea}

COLoRA is built around one simple idea: in the joint-object regime, a miss should not automatically force the request to wait for GPU promotion. Instead, the runtime should separate the hot-set problem from the cold-miss problem. The hot set should remain on the GPU because reuse justifies residency. The cold set should remain executable from CPU memory so that an individual miss can be handled immediately, even if the object is not worth promoting yet.

At a high level, COLoRA combines three mechanisms. First, it uses an asymmetric memory hierarchy with a bounded GPU hot cache and a CPU-resident full replica. Second, it uses CPU-first miss handling for the current request and decouples promotion into an asynchronous, selective background action. Third, it overlaps GPU hits with CPU misses whenever both appear in the same dispatch, so the fallback path can hide behind useful GPU work instead of serializing it.

\subsection{Asymmetric Joint-Object Memory Hierarchy}
\label{subsec:memory-hierarchy}

COLoRA manages LoRA residual weights with an asymmetric memory hierarchy. GPU memory is treated as a scarce hot cache that should hold only reusable joint objects. CPU memory holds the complete read-only source pool for all joint objects and therefore acts as both the cold layer and the execution source for fallback. This asymmetry is important because it removes write-back complexity from eviction and ensures that an evicted object remains immediately executable from host memory.

The GPU cache is populated selectively rather than eagerly. Each joint object maintains lightweight utility state derived from recent accesses, and promotion decisions are filtered through minimum-hit thresholds, bounded promotion queues, and cooldown windows. Intuitively, COLoRA promotes objects that show a plausible hot rebound, while suppressing repeated queue traffic for objects that remain cold. This policy is designed to preserve cache precision under the enlarged key space without letting promotion traffic dominate the worker.

\subsection{CPU-First Non-Blocking Miss Handling}
\label{subsec:cpu-first}

On the execution path, COLoRA distinguishes between the dense expert backbone and the sparse LoRA residual. The dense base expert computation remains on GPU. Only the expert-specific LoRA residual is allowed to use the heterogeneous path. This split is crucial because it limits the scope of fallback to the incremental computation that caused the residency problem in the first place, rather than moving the entire MoE layer off GPU.

When a joint object hits in the GPU cache, COLoRA serves it through the regular GPU LoRA path. When it misses, the current request takes a CPU fallback path by default instead of blocking on promotion. The runtime may still enqueue the missed object for promotion, but promotion is treated as an optimization for future requests, not a prerequisite for the current one. This is the central paradigm shift of the design: the system no longer interprets every miss as ``move weight first, execute later.'' It interprets a miss as ``execute now from the cold layer, and only promote if reuse signals justify it.''

\subsection{Low-Overhead CPU Fallback for Sparse Decode Misses}
\label{subsec:cpu-fallback}

The CPU fallback path is explicitly engineered for the small and irregular shapes that appear in decode. Instead of relying only on generic framework operators, COLoRA groups miss work by adapter, uses MoE-specific CPU kernels when available, and keeps the per-miss control path compact. The implementation already exposes AVX-512 BF16 kernels specialized for the gate, up, and down phases of MoE LoRA execution, which matches the paper's requirement that CPU fallback be a viable latency path rather than a correctness-only fallback.

This design choice should be framed carefully in the paper. The claim is not that CPU is generally faster than GPU. The claim is that for cold decode misses with poor reuse, the fixed cost of blocking promotion can exceed the fixed cost of executing the residual directly on CPU. COLoRA wins only if the fallback path stays cheap enough to make that trade favorable on the tail.

\subsection{Overlapped Heterogeneous Execution}
\label{subsec:overlap}

The final piece is overlap. Mixed batches often contain both GPU-cache hits and misses. COLoRA exploits this structure by launching CPU miss work asynchronously when useful GPU hit work is available, then joining the miss result only at the point where the residual output must be merged. This organization lets CPU compute, GPU compute, and host--device movement proceed concurrently when the batch structure allows it.

The runtime also keeps explicit control signals for this path, including promotion queue depth, CPU queue depth, queueing time, overlap ratio, fallback degradations, and promotion-drop causes. These counters matter for an ASPLOS paper because they make the system explainable under load. If COLoRA helps, the paper can attribute the benefit to fewer blocking promotions and better overlap. If it stops helping, the same counters can reveal whether the cause is CPU queue buildup, insufficient GPU hit work, or a pathological cold burst.

\subsection{Expected Benefit and Graceful Degradation}
\label{subsec:graceful-degradation}

COLoRA is not designed to force heterogeneity into every workload. If locality is strong and the active joint objects fit comfortably in GPU memory, the system should naturally behave close to a GPU-only runtime because most requests hit the hot cache and the CPU path rarely triggers. The design matters precisely in the fragmented regime identified by the case study, where cold misses are common enough that a stable miss-time fallback is more valuable than ever more aggressive blocking promotion.


\section{Scope and Limitations}
\label{sec:limitations}

COLoRA assumes that CPU memory can hold the complete joint-object source pool. This assumption should be stated explicitly because the design depends on a CPU-resident full replica as the cold layer.

COLoRA is most attractive when the workload exhibits a hot--cold split: a reusable GPU-resident head plus a sparse cold tail. If the cold fraction becomes too large, CPU queueing can erode the benefit. Likewise, COLoRA reduces exposure to blocking promotion, but it does not eliminate host--device transfer costs or hardware sensitivity in the CPU path. NUMA placement, CPU vector support, and contention still matter.

\section{Evaluation}
\label{sec:evaluation}

We evaluate COLoRA along five dimensions that correspond directly to the claims made in this paper. In particular, we ask the following questions:

\begin{itemize}
\item \textbf{RQ1: End-to-end latency.} Does COLoRA reduce end-to-end decode latency under multi-LoRA workloads?
\item \textbf{RQ2: Tail latency.} How much does COLoRA improve P95/P99 latency under long-tail expert--adapter access patterns?
\item \textbf{RQ3: Miss handling.} How effective is COLoRA's non-blocking miss path, which combines a GPU hot cache with CPU fallback, relative to conventional load-then-run execution?
\item \textbf{RQ4: Component contribution.} What performance gain is attributable to each major component, including CPU fallback and speculative dispatch?
\item \textbf{RQ5: System behavior.} How does COLoRA affect cache hit rate, miss cost, overlap efficiency, and overall resource utilization?
\end{itemize}

\subsection{Hardware and System Configuration}
\label{Configuration}
All experiments are conducted on a single server with the following configuration:

\begin{itemize}
\item \textbf{GPU:} A100 80GB / RTX 6000
\item \textbf{CPU:} Intel Xeon Gold 5520+ (1 socket, 28 physical cores, 56 hardware threads, AVX-512 support)
\item \textbf{Memory:} 376GB RAM
\item \textbf{Interconnect:} PCIe Gen4 / NVLink
\end{itemize}

Unless otherwise noted, the CPU fallback path uses AVX-512-optimized kernels.


\subsection{Models and Workloads}
\label{subsec:workloads}
We instantiate COLoRA on the public \textbf{Qwen3-VL-30B-A3B-Instruct} checkpoint. Although Qwen3-VL is a multimodal model, our evaluation targets its MoE text decoder because COLoRA operates on decode-time expert--adapter management. The evaluated decoder has 48 layers, hidden size 2048, 128 routed experts per layer, and top-8 routing for each token at every sparse layer. For adapter experiments, we use rank-16 LoRA modules applied to the attention and MLP projections, including the MoE gate, up, and down projections.

All workloads target the decode worker in a PD-separated serving pipeline. We collect real decode-phase routing traces by replaying a fixed corpus of 512 ShareGPT V3 requests through the model and recording the router-selected experts at each token. We then pair this routing stream with semi-realistic adapter-invocation traces derived from the Azure Functions workload trace, which preserves realistic tenant-arrival skew while allowing controlled variation in adapter-side multiplexing.

Unless otherwise specified, we evaluate both independent and correlated tenant-to-adapter mappings and sweep the active adapter cardinality over $\{8, 32, 128\}$. We further vary the decode batch size (typically 1--4) and the request concurrency at the decode worker to test how COLoRA behaves as joint expert--adapter fragmentation and system contention increase together.

\subsection{Baselines}
\label{subsec:baselines}

The main comparison should include an optimized GPU-only baseline that keeps GPU as the only execution target and uses the best available caching and prefetching policy. A second baseline should adopt joint-object cache management but still resolve misses by blocking promotion, which isolates the value of heterogeneous miss handling from the value of finer-grained residency alone. Finally, the paper should include COLoRA ablations, including at least variants without CPU fallback, without asynchronous overlap, and without the optimized MoE CPU kernel path.

% We compare COLoRA with the following baselines:

% GPU-Only:
% All LoRA–expert weights reside on GPU. No offloading.

% Naive Offload (Load-Then-Run):
% On a cache miss, weights are synchronously transferred from CPU to GPU before execution.

% [Optional] Static Partitioning:
% A static subset of LoRA–expert pairs is pinned in GPU memory.

\subsection{End-to-End Latency}
\label{subsec:e2e-lat}

% The primary latency metric should be TPOT at mean, P90, and P99 because the paper is fundamentally about decode stability. Throughput is also necessary to show that tail-latency gains are not purchased with unacceptable aggregate slowdown. TTFT should be reported as well to verify that decode-side changes do not create obvious collateral damage in the broader serving path.

We first evaluate whether COLoRA improves end-to-end decoding latency.
% TODO:
Figure X shows the average and tail latency (P95/P99) across different workloads.

COLoRA reduces average latency by up to X\% and P99 latency by up to Y\% compared to Naive Offload.

This improvement comes from avoiding blocking weight transfers on cache misses.

Instead of stalling on load-then-run, COLoRA executes miss requests via CPU fallback,
thereby converting latency spikes into more stable execution paths.

\subsection{Tail Latency Reduction}
\label{subsec:tail-lat}

% The evaluation should also expose the runtime mechanics that explain the end-to-end behavior. At minimum, the paper should report cache hit rate and miss breakdown, CPU queue depth and queue wait, CPU compute time, GPU compute time, overlap ratio, promotion-drop breakdown, and host memory footprint. These measurements let the paper demonstrate whether COLoRA improves P99 because it changes miss handling in the intended way, rather than because of an unrelated system effect.

We analyze the impact of COLoRA on tail latency under skewed access patterns.

Under high LoRA concurrency, Naive Offload exhibits significant tail amplification,
with P99 latency increasing sharply.

COLoRA reduces P99 latency by up to X\times.

This is because COLoRA decouples cache miss from request blocking.

Long-tail accesses no longer incur full transfer latency,
but are instead handled via CPU fallback with bounded overhead.

% 👉 关键词要反复打：

% “decoupling miss from blocking”

% “bounded miss cost”

\subsection{Effectiveness of Non-Blocking Miss Handling}
\label{subsec:miss-handling}
We evaluate how COLoRA’s miss handling policy affects system behavior.

Metrics need to report:

miss rate

miss latency

fallback ratio

COLoRA significantly reduces effective miss penalty.

While the miss rate remains similar, the average miss cost decreases by X\%,
as many misses are served via CPU fallback instead of blocking transfers.

This demonstrates that optimizing miss handling policy is as important as improving cache hit rate.

Even when misses are unavoidable, COLoRA ensures they are handled in a non-blocking manner.

\subsection{Breakdown Analysis}
\label{subsec:breakdown}
We break down end-to-end latency into major components.

TODO: figures

stacked bar:

GPU compute

CPU fallback

transfer

idle/wait

COLoRA reduces idle time caused by blocking transfers.

Although CPU fallback introduces additional compute,
it overlaps with GPU execution and avoids long stalls.

\subsection{Ablation Study}
\label{subsec:ablation}
We evaluate the contribution of each component.

- Full COLoRA
- w/o fallback (always load-then-run)
- w/o speculative dispatch

Removing fallback significantly increases tail latency,
confirming its role in mitigating miss-induced stalls.

Speculative dispatch further improves performance by overlapping CPU and GPU execution.

% “救命 checklist”

% 要保证：

% ✅ 有 end-to-end latency 图
% ✅ 有 P99 改善
% ✅ 有 baseline 对比
% ✅ 有 1–2 个 insight（不是纯图）

% ============================ GPT Notes ============================
% Compact outline:
% - The paper targets PD-separated decode workers that serve an MoE base
%   model together with many concurrently active LoRA adapters.
% - The core observation is that the effective managed object is no longer
%   an expert or an adapter alone, but a joint expert-adapter object whose
%   locality is much weaker and more fragmented.
% - A trace-driven case study shows a consistent chain: flatter popularity,
%   worse temporal locality, higher miss rates under the same budget, and a
%   tail penalty that turns on early once the system leaves the single-LoRA
%   regime.
% - The resulting systems problem is not just cache sizing. It is how to
%   handle unavoidable misses without exposing weight migration on the decode
%   critical path.
% - COLoRA addresses this problem with an asymmetric memory hierarchy,
%   CPU-first non-blocking miss handling, selective asynchronous promotion,
%   and overlapped heterogeneous execution.
% - The evaluation should focus on P99 TPOT, throughput, miss-handling
%   breakdown, and the cost profile of the CPU fallback path.
%
% Reverse outline:
% - Introduction thesis: joint-object fragmentation makes decode misses
%   structural, so the system must optimize miss handling, not just miss
%   avoidance.
% - Introduction paragraph 1: PD-separated MoE multi-LoRA decode serving is
%   tail-sensitive.
% - Introduction paragraph 2: the managed object becomes a joint
%   expert-adapter object.
% - Introduction paragraph 3: the case study shows flatter popularity, worse
%   locality, and early tail turn-on.
% - Introduction paragraph 4: GPU-only load-then-run is a poor default once
%   cold misses become structural.
% - Introduction paragraph 5: COLoRA handles hot and cold objects
%   differently through a dual-path runtime.
% - Background thesis: the joint-object view changes both the memory problem
%   and the execution problem.
% - Challenges thesis: the runtime must control residency churn, CPU fallback
%   cost, and heterogeneous synchronization together.
% - Design thesis: COLoRA keeps hot objects on GPU, executes cold misses from
%   CPU, and only promotes when reuse signals justify it.
%
% Self-review checklist:
% - Clarity: Each section has an explicit thesis sentence and stable
%   terminology around joint expert-adapter object, GPU hot cache, and CPU
%   fallback.
% - Flow: The story moves from setting -> failure mode -> evidence -> design
%   question -> design overview.
% - Terminology consistency: The paper reasons about logical joint objects
%   and notes that the implementation realizes them as projection-specific
%   entries.
% - Unsupported claims: The draft avoids claiming measured end-to-end gains
%   that are not yet written down.
% - Missing evidence: The final paper still needs full end-to-end COLoRA
%   results, especially TPOT/P99 comparisons against strong GPU-only
%   baselines and ablations for queue control and CPU kernels.
%
% Claim-evidence map:
% - Claim: In MoE multi-LoRA decode serving, the effective managed object is
%   a joint expert-adapter object rather than an expert or adapter alone.
%   Evidence: The request needs LoRA residuals only for the experts selected
%   by the router; the implementation also keys GPU residency by projection,
%   adapter, layer, and expert. Status: supported.
% - Claim: Joint-object modeling significantly fragments popularity and
%   weakens temporal locality. Evidence: Top-1000 access coverage drops from
%   89.1% to 40.7% and 32.6%; finite P99 reuse distance grows from 1,393 to
%   2,041 and 4,521; cold first touches rise from 1.0 to 9.7 and 13.6 per
%   10k events. Status: supported.
% - Claim: Under the same cache budget, joint-object fragmentation amplifies
%   misses enough to become a tail-latency problem. Evidence: At 2048
%   objects, miss rate rises from 0.00159 in B0 to 0.01095 and 0.01476 in
%   B1/B2; the 8192-object sweep shows a P99 latency-proxy penalty of
%   411.7 ms by four LoRAs and a broad high plateau afterward. Status:
%   supported.
% - Claim: Tail requests are dominated by cold joint-object touches rather
%   than by a uniform slowdown across all requests. Evidence: At the
%   65536-object cold-miss floor, average requests touch 83.6 and 116.8 cold
%   joint objects, while P95+ tail requests touch 912.8 and 1154.0. Status:
%   supported.
% - Claim: COLoRA's key design shift is to treat a miss as an execution-path
%   decision instead of a mandatory blocking promotion event. Evidence: The
%   runtime maintains a GPU hot cache plus a CPU full replica; the configured
%   default miss policy is cpu_first; misses can execute immediately on CPU
%   while promotion is queued asynchronously for future reuse. Status:
%   supported.
% - Claim: CPU fallback is plausible only if the runtime keeps the fixed cost
%   of sparse decode misses low. Evidence: The implementation includes
%   MoE-specific CPU kernels, adapter-grouped miss execution, asynchronous CPU
%   queueing, and explicit queue/overlap counters. Status: supported for
%   mechanism, needs end-to-end performance evidence for effect.
% - Claim: Overlap between GPU hits and CPU misses is an important part of
%   the design rather than an implementation detail. Evidence: The hybrid
%   path tracks CPU compute time, GPU compute time, queue wait, overlap
%   ratio, and fallback degradations; mixed hit-miss tests exercise this
%   path. Status: supported for mechanism, needs end-to-end performance
%   evidence for effect.



\section{Related Work}
\section{Conclusion}

\bibliographystyle{plain}
\bibliography{COLoRA-reference}

\end{document}
\endinput
%%
%% End of file `sample-sigconf.tex'.
